% Coupling paper, first complete draft (27 Oct 2026).
% Build (offline): pdflatex main; bibtex main; pdflatex main; pdflatex main.
% The Springer Nature template (sn-jnl.cls, Biomechanics and Modeling in
% Mechanobiology) is not installed here; the manuscript uses the article class
% with natbib (author-year), 12 pt, single column and line numbers. The Springer
% template will be applied at submission.
\documentclass[11pt,a4paper]{article}
\usepackage[utf8]{inputenc}
\usepackage[T1]{fontenc}
\usepackage{lmodern}
\usepackage[a4paper,margin=25mm]{geometry}
\usepackage{amsmath,amssymb,bm,mathtools}
\usepackage{graphicx}
\usepackage{booktabs}
\usepackage{array}
\usepackage{caption}
\usepackage{subcaption}
\usepackage[authoryear,round]{natbib}
\usepackage{algorithm}
\usepackage{algpseudocode}
\usepackage{lineno}
\usepackage{xcolor}
\usepackage[colorinlistoftodos,textsize=footnotesize]{todonotes}
\usepackage[hidelinks]{hyperref}
\usepackage{microtype}
\setlength{\emergencystretch}{2em}
\renewcommand{\topfraction}{0.9}\renewcommand{\bottomfraction}{0.6}\renewcommand{\textfraction}{0.05}\renewcommand{\floatpagefraction}{0.8}
\setcounter{topnumber}{3}\setcounter{totalnumber}{4}
\captionsetup{font=small,labelfont=bf}
\linenumbers

% Paths of the figures (not copied into this folder).
\newcommand{\assets}{../../../assets/}
\newcommand{\figs}{../}
% Figure in preparation by another script: box with the file name as a fallback.
\newcommand{\figorbox}[2]{\IfFileExists{#1}{\includegraphics[#2]{#1}}{\fbox{\parbox[c][5cm][c]{0.9\linewidth}{\centering figure file \texttt{\detokenize{#1}} not yet generated}}}}

% Notation (Klempt et al. 2024, 2026).
\newcommand{\phiD}{\phi_{3D}}
\newcommand{\lv}{\textstyle\sum_i\bar\phi_i}
\newcommand{\nc}{\mathbf n_c}
\newcommand{\fterm}{\ensuremath{-r\tfrac{c}{k+c}\,\nc\cdot\nabla\phiD}}
\newcommand{\TODO}[1]{\todo[inline]{TODO: #1}}

\title{Coupling a multi-species point model to a continuum biofilm growth model: verification in three codes and the effect of the composition on growth and stress}
\author{Authors to be agreed\thanks{Author list, order and affiliations to be agreed by all contributors before the manuscript is shared; corresponding author to be named.}}
\date{Draft of 27 October 2026}

\begin{document}
\maketitle

\begin{abstract}
Continuum models of biofilm growth predict where biomass forms and the stresses it causes, but treat the biofilm as one material. Multi-species point models describe the composition, the living fraction and the response to an antibiotic, but have no space. This work couples the multi-species point model of Klempt et al.\ (2026), used as published, to the Hamilton-principle growth model of Klempt et al.\ (2024) at every integration point, without changing the equations of either model. The point model receives the local nutrient and antibiotic from the fields, and its living volume fraction $\sum_i\bar\phi_i$ sets the consumption $g\,\phi_{3D}(\sum_i\bar\phi_i)\,c$ of both. The structure of the coupled equations fixes what the composition can change: it reaches growth and stress only through the nutrient and the front speed of the growth equation, because the growth law counts all biomass, so the growth is bounded by $k_\alpha T^*$ and the front speed by a factor of two. A model in which the composition is to change the mechanics therefore needs a growth rate or a front speed tied to the living biomass or to the species, and the coupling presented here is the base on which such a term can be added and measured. The quasi-steady nutrient depends on the composition through one product and is depleted over the half edge of the domain (Thiele modulus 1). In the point model the antibiotic acts through a balance whose decay rate grows as the inverse square of the species' volume fraction, so a dense species is removed slowly and a thinly present one at once. Examples with two species and an antibiotic, computed in a finite-volume code and in user elements of ANSYS and Abaqus and verified by a manufactured solution and grid convergence indices, confirm this: the antibiotic removes most of the living biomass of the seed (0.96 to 0.15 when the species start apart, 0.96 to 0.04, ten times faster, when mixed), while the growth changes by about 1\,\% and the converged seed stress within its grid uncertainty; the local nutrient decides where the antibiotic acts; and the state that newly reached material receives, which neither model fixes, decides in a maze whether the biofilm recovers. The same holds around a human molar crown. A second-order treatment of the front term is needed for the stress.
\end{abstract}

\noindent\textbf{Keywords:} biofilm growth; multi-species model; antibiotic; extended Hamilton principle; code verification; grid convergence index; growth-induced stress

\section{Introduction}
\label{sec:intro}

Oral biofilms on teeth and implants are communities of several bacterial species. Their composition changes over days: early colonisers prepare the surface for later ones, and pathogenic species appear only when the community allows it \citep{Heine2025PeriImplant}. Treatments such as antibiotics act on the composition, while the mechanical load on the surface comes from the amount and the growth of the biomass. A model that predicts the stresses for a given treatment therefore has to describe both: where the biofilm grows, and what it is made of.

Two kinds of model exist for this. Continuum growth models describe the amount of biomass in space. \citet{WannerGujer1986} introduced the one-dimensional multispecies biofilm, \citet{Eberl2001Biofilm} a density-dependent diffusion model of the spreading. \citet{Feng2021SymbioticBiofilm} model a two-species oral biofilm in which the spreading follows the summed growth of the species, and \citet{Soleimani2023CoAggregation} the co-aggregation of two species in a user element. \citet{Klempt2024DiffusionDrivenGrowth} derive the growth of a biofilm from the extended Hamilton principle \citep{JunkerBalzani2021ExtendedHamilton}, with a phase field $\phiD$, a nutrient $c$ and the growth variable $\alpha$ of a multiplicative split $\mathbf F=\mathbf F_e\mathbf F_g$, and compute the resulting stresses; their model has one biomass and no composition. Point models describe the composition of a community at one point. \citet{Klempt2026ContinuumBacterialGrowth} derive a multi-species model from the same variational principle, with volume fractions $\phi_i$ and living fractions $\psi_i$ of $N$ species, an interaction matrix $\mathbf A$ and the response of each species to an antibiotic; \citet{Fritsch2025BayesianMicrofilms} calibrate its parameters by Bayesian updating. These models have no space: the nutrient $c^*$ and the antibiotic $\alpha^*$ are given inputs.

The point model needs local values of $c^*$ and $\alpha^*$, which a growth field can supply. The growth field needs to know how much of the biomass is alive, because only living biomass consumes nutrient. Connecting the two at every material point is the natural step, but it raises questions that neither paper answers: which variables are passed, whether either model has to be changed, which new assumptions are needed (for example the composition of biofilm that a point reaches later), and whether the composition reaches the growth and the stress at all.

This work couples the point model of \citet{Klempt2026ContinuumBacterialGrowth} to the growth model of \citet{Klempt2024DiffusionDrivenGrowth}. To our knowledge it is the first coupling of a variational multi-species point model to a continuum growth model with stress, and it is made without changing either set of equations: the point model runs at every integration point as published, receives $c^*=c^*_0c$ and $\alpha^*=\alpha^*_0a$ from the nutrient and antibiotic fields and returns the living volume fraction $\lv$, which sets the consumption $g\,\phiD(\lv)\,c$ in the field equations. The contributions are the following.
\begin{enumerate}
\item The main result is a property of the coupled equations. In the model of \citet{Klempt2024DiffusionDrivenGrowth} the composition reaches the growth and the stress along one path only, through the nutrient and the front speed of the growth equation, because the growth law counts all biomass; the growth is bounded by $k_\alpha T^*$ and the front speed changes by at most a factor of two (Section~\ref{sec:structure}). The examples confirm it: the antibiotic removes most of the living biomass of the seed and changes the growth by about $1\,\%$ and the converged seed stress within its grid uncertainty. Any model in which the composition is to change the mechanics therefore needs a growth rate or a front speed tied to the living biomass or to the species; the coupling presented here is the base on which such a term can be added, and the size of its effect is estimated from the runs (Section~\ref{sec:discussion}).
\item Three further properties are derived from the published equations and read against the examples: the quasi-steady nutrient with its Thiele modulus, the balance by which the antibiotic acts in the point model, whose decay rate grows as the inverse square of the volume fraction of a species, and the time link, which is equivalent to a scaling of the viscosities and is therefore fixed by the data.
\item The coupling is implemented independently in a finite-volume code and in user elements of ANSYS and Abaqus, and verified by replay of the point model, a manufactured solution and grid convergence indices. The front term \fterm{} of the growth equation sets the accuracy of the stress and needs a second-order treatment.
\item The state that newly reached material receives is fixed by neither model. It is stated as an open modelling choice, and the examples show its effect: in a maze it decides whether the biofilm recovers after the antibiotic.
\end{enumerate}
The examples use two species and an antibiotic, with the synthetic data of \citet{Fritsch2025BayesianMicrofilms} and the constants of \citet{Klempt2024DiffusionDrivenGrowth,Klempt2026ContinuumBacterialGrowth}, on a cube, against obstacles, in a maze and around a human molar crown.
Section~\ref{sec:models} states the models, the coupling and what its structure implies, Section~\ref{sec:verification} the implementations and their verification in short (details in the appendix), Section~\ref{sec:results} the examples, Section~\ref{sec:discussion} the limits of the coupling, and Section~\ref{sec:conclusions} concludes.

\section{Models and coupling}
\label{sec:models}

\subsection{Growth field}
\label{sec:field}

The biofilm is described by the phase field $\phiD\in[0,1]$, the nutrient $c$ and the growth variable $\alpha$ of \citet{Klempt2024DiffusionDrivenGrowth}, with $\mathbf F=\mathbf F_e\mathbf F_g$, $\mathbf F_g=\alpha\mathbf I$ and $\alpha(0)=1$, so that the growth is $\alpha-1$. After division by the viscosities, their Eqs.~34 to 36 read
\begin{align}
\dot\phiD &= \beta\,\nabla^2\phiD + k_\alpha\,\alpha - r\,\frac{c}{k+c}\,\nc\cdot\nabla\phiD, \qquad \nc=\frac{\nabla c}{|\nabla c|}, \label{eq:34}\\
\dot c &= d\,\nabla^2 c - g\,\phiD\Bigl(\lv\Bigr)\,c, \label{eq:35}\\
\dot\alpha &= k_\alpha\,\phiD, \label{eq:36}
\end{align}
with zero flux on the boundary except where the nutrient is held. The last term of Eq.~\eqref{eq:34} moves the edge of the biofilm up the nutrient gradient with the speed $r\,c/(k+c)$; we call it the front term. Two details differ from the equations as printed, both confirmed by the authors: the consumption of Eq.~\eqref{eq:35} is $g\,\phi\,c$ (the printed $g\,\phi$ is a misprint), and the front direction is regularised as $\nabla c/(|\nabla c|^2+\varepsilon^2)^{1/2}$ with $\varepsilon^2=10^{-6}$ ($\varepsilon^2=10^{-8}$ changes the seed stress by $0.1\,\%$). The factor $\lv$ in Eq.~\eqref{eq:35} is the coupling term of this work (Section~\ref{sec:coupling}); in \citet{Klempt2024DiffusionDrivenGrowth} the whole biofilm consumes. The mechanical term $\phi\,\mu\,(\mathbf I:\mathbf C_e-3)$ of their Eq.~30, which they neglect because $\mu$ is small against the other parameters, is neglected here as well, so the mechanics is coupled one way. The elastic response is neo-Hookean with the stiffness $E(\phiD)=E\,(\phiD^2+f)$, $E=10$~Pa, $\nu=0.49$ and a void floor $f=10^{-3}$ (an assumption that keeps the empty cells solvable; $f\to0$ changes the seed stress by $0.13\,\%$).

\subsection{Point model}
\label{sec:point}

At a material point the multi-species model of \citet{Klempt2026ContinuumBacterialGrowth} describes the volume fractions $\phi_i$ and the living fractions $\psi_i$ of $N$ species and the void $\phi_0$, with $\sum_i\phi_i+\phi_0=1$. It follows from the extended Hamilton principle with the free energy $\Psi=-\tfrac12c^*\,\bar{\bm\phi}^{\mathrm T}\mathbf A\,\bar{\bm\phi}$, $\bar\phi_i=\phi_i\psi_i$, a dissipation with the viscosities $\eta_i$, $\eta_{\phi,i}$ and the constraint enforced by the multiplier $\gamma$:
\begin{equation}
\begin{aligned}
(\eta_{\phi,i}+\eta_i\psi_i^2)\,\dot\phi_i+\eta_i\phi_i\psi_i\,\dot\psi_i
&= c^*\psi_i\textstyle\sum_j A_{ij}\bar\phi_j-\gamma-(\text{barrier}),\\
\eta_i\,(\phi_i\psi_i\,\dot\phi_i+\phi_i^2\,\dot\psi_i)
&= c^*\phi_i\textstyle\sum_j A_{ij}\bar\phi_j-\alpha^* b_i\,\psi_i-(\text{barrier}),
\end{aligned}
\label{eq:point}
\end{equation}
where $A_{ij}$ are the interactions, $b_i$ the sensitivity of species $i$ to an antibiotic of concentration $\alpha^*$, $c^*$ the nutrient, and the barrier terms keep $\phi_i,\psi_i$ in $(0,1)$. The antibiotic lowers $\psi_i$ through the term $\alpha^*b_i\psi_i$, against the nutrient term $c^*\phi_i\sum_jA_{ij}\bar\phi_j$ (Section~\ref{sec:structure}). Our implementation reproduces the fourteen numerical examples of the paper (two printed details read differently: Eq.~17 as printed carries a $+\gamma$ that the figures do not show, and the four-species interaction matrix of Eq.~19 reproduces the figures without its prefactor $\tfrac12$) and the figures of \citet{Fritsch2025BayesianMicrofilms} within $0.013$. The model is used as published, with its own constraint and its implicit time integration (steps of $10^{-4}$ in its own time unit, Newton's method). For the examples the synthetic two-species data of \citet{Fritsch2025BayesianMicrofilms}, case~I, are used with their true parameters: $\mathbf A=\begin{psmallmatrix}1&0.1\\0.1&1\end{psmallmatrix}$, $b=(1,2)$, $\eta=(1,2)$, $c^*_0=100$.

\subsection{Coupling}
\label{sec:coupling}

The two models are connected at every integration point without changing either (Figure~\ref{fig:schematic}, Algorithm~\ref{alg:coupled}).
\begin{enumerate}
\item \emph{Fields to point model.} The point model sees $c^*=c^*_0\,c$ and $\alpha^*=\alpha^*_0\,a$, with $a$ the antibiotic, transported by a second field with the diffusivity and uptake of the nutrient, $\dot a=d\,\nabla^2a-g\,\phiD(\lv)\,a$. Between the increments $t_n$ and $t_{n+1}$ of the field both are interpolated linearly in time over the inner steps of the point model,
\begin{equation}
c^*(t)=c^*_0\Bigl[\tfrac{t_{n+1}-t}{\Delta t}\,c_n+\tfrac{t-t_n}{\Delta t}\,c_{n+1}\Bigr],\qquad t_n\le t\le t_{n+1}.
\label{eq:cstar}
\end{equation}
\citet{Klempt2026ContinuumBacterialGrowth} already use a time-dependent $c^*$ (their four-species case~3), so the point model accepts such an input without change.
\item \emph{Point model to fields.} Only the living biomass consumes: the uptake of nutrient and antibiotic is $g\,\phiD(\lv)$, Eq.~\eqref{eq:35}. The point model normalises its volume, so $\bar\phi_i\in[0,1]$ is the living volume fraction of species $i$ at the point and $\phiD\lv\le\phiD$ the living volume fraction in the field; the void $\phi_0$ has no living fraction and is not part of the sum. A dead biofilm consumes nothing. The growth law~\eqref{eq:36} is unchanged: $\alpha$ measures the biomass produced, alive or dead.
\item \emph{Time.} \citet{Klempt2024DiffusionDrivenGrowth} normalise time with a reference time chosen for the problem at hand. We take $T^*=1$ as the duration covered by the point model's data, so that both models run on the same time: the point model advances $s\,\Delta T^*$ per increment, with $s$ the length of that window in its own units, not a fitted value ($s=0.1$ for case~I, 1000 steps of $10^{-4}$).
\item \emph{Start of the point model.} It starts at a point once $\phiD$ there reaches $\phi_{\min}=0.001$. Where a seed is placed it starts from the seed composition. A point that the spreading biofilm reaches later takes the living composition of its colonised neighbours (their amounts divided in the ratio of $\bar\phi_i$) and starts alive ($\psi_i=0.999$, $\gamma=0$), because the biofilm spreads through the growth of its living cells. The papers do not fix this rule; it is an assumption, and the alternative of starting from the seed composition is compared in Section~\ref{sec:newpoints}.
\end{enumerate}
What this coupling can and cannot change follows from the structure of the equations and is stated next.

\begin{figure}[tbp]
  \centering
  \includegraphics[width=0.62\linewidth]{\assets fig_coupling_schematic_agreed.png}
  \caption{The coupling: the fields $\phiD$, $c$ and $a$ of the growth model, the point model at every integration point with the inputs $c^*=c^*_0c$, $\alpha^*=\alpha^*_0a$ and the output $\lv$, and the mechanics with $\mathbf F_g=\alpha\mathbf I$.}
  \label{fig:schematic}
\end{figure}

\begin{algorithm}[tbp]
\caption{One increment $t_n\to t_{n+1}$ of the coupled run (user-element form)}
\label{alg:coupled}
\begin{algorithmic}[1]
\Require at every integration point $\alpha_n$, the point-model state $(\phi_i,\psi_i,\gamma)_n$ and the pair $(c,a)$ seen in the previous increment; the fields $\phiD$, $c$, $a$ of $t_n$
\Ensure displacements, stresses and state of $t_{n+1}$
\Repeat \Comment{Newton iteration of the equilibrium}
  \ForAll{integration points}
    \State $\alpha_{n+1}\leftarrow\alpha_n+k_\alpha\,\phiD\,\Delta t$ \Comment{Eq.~\eqref{eq:36}}
    \State $\mathbf F_e\leftarrow\mathbf F\,\mathbf F_g^{-1}$, $\mathbf F_g=\alpha_{n+1}\mathbf I$; stress and tangent from $W(\mathbf F_e)$ and $E(\phiD)$
    \If{$\phiD\ge\phi_{\min}$ or the point model has started}
      \If{not yet advanced in this increment}
        \State $(c_n,c_{n+1})$, $(a_n,a_{n+1})$ from the previous and the current value at the point
        \State $n_{\rm sub}$ steps of $s\Delta t/n_{\rm sub}$ of Eq.~\eqref{eq:point} with $c^*$, $\alpha^*$ from Eq.~\eqref{eq:cstar}
        \State store $(\phi_i,\psi_i,\gamma)_{n+1}$ and $\lv$
      \EndIf
    \EndIf
  \EndFor
\Until{the residual meets the convergence criterion}
\State advance $\phiD$ (Eq.~\eqref{eq:34}), $c$ and $a$ (Eq.~\eqref{eq:35} with $\lv$)
\end{algorithmic}
\end{algorithm}

\subsection{What the structure of the coupled equations implies}
\label{sec:structure}

The main result of this work follows from Eqs.~\eqref{eq:34} to \eqref{eq:point} alone: the composition reaches the growth and the stress along a single path, which bounds what any change of the composition can do to the mechanics, and which says what a model needs if the composition is to change the stress. Three further properties follow from the same equations; all four are used to read the examples.

\paragraph{The path from the composition to the stress.} The composition enters the field equations only through the factor $\lv$ in the consumption of Eq.~\eqref{eq:35}. The nutrient enters Eq.~\eqref{eq:34} only in the front term, through the speed $r\,c/(k+c)$ and the direction $\nc$, and it does not enter the growth law~\eqref{eq:36} at all. The composition therefore reaches $\phiD$, $\alpha$ and the stress along one path only,
\begin{equation}
\lv\ \xrightarrow{\ \eqref{eq:35}\ }\ c\ \xrightarrow{\ \text{front term of \eqref{eq:34}}\ }\ \phiD\ \xrightarrow{\ \eqref{eq:36}\ }\ \alpha\ \longrightarrow\ \boldsymbol\sigma,
\label{eq:path}
\end{equation}
and it has no effect wherever that path is cut: in the interior of a biofilm, where $\nabla\phiD=\mathbf 0$, and wherever the nutrient is uniform, where $\nabla c=\mathbf 0$ and the front term vanishes. Two bounds follow. Since $0\le\phiD\le1$, integrating Eq.~\eqref{eq:36} gives $\alpha(T)-1=k_\alpha\int_0^T\phiD\,dt\le k_\alpha T$; in the seed of the examples $\alpha-1$ reaches $0.91$ of this bound at $T^*=1$, so no change of the composition can raise the seed growth by more than about $10\,\%$, and it can lower it only through the edge of the seed, where $\phiD<1$ enters the mean. And the front speed $r\,c/(k+c)$ with $k=1$ and $0\le c\le1$ lies between $0$ and $r/2$, so the composition changes the speed of the edge by at most a factor of two through the nutrient, and in practice by the ratio $(c/(1+c))$ at the two nutrient levels (by $37\,\%$ between the seed means $0.54$ and $0.91$ of Section~\ref{sec:composition}). The stress is set by the jump of $\alpha$ across the edge of the seed, so the composition acts on the stress through the position and sharpness of that edge, not through the amount of growth. Both bounds are properties of the published equations, and they have a consequence for the modelling: a composition that is to change the stress needs a second path, a growth law in which the rate counts the living biomass, $\dot\alpha=k_\alpha\phiD\lv$, a front speed multiplied by $\lv$, or a stiffness that depends on the species, and each of these is a change of the model (Section~\ref{sec:discussion}).

\paragraph{The nutrient is quasi-steady and depends on the composition through one product.} With the constants of Table~\ref{tab:parameters} the nutrient adjusts within about $4\times10^{-8}\,T^*$, so Eq.~\eqref{eq:35} reduces to the linear problem
\begin{equation}
\nabla^2 c=\frac{\Lambda^2}{\ell^2}\,\kappa\,c,\qquad \kappa=\phiD\lv,\qquad \Lambda^2=\frac{g\ell^2}{d}=1,
\label{eq:thiele}
\end{equation}
with $\ell=1$~mm the half edge of the cube and $\Lambda$ the Thiele modulus of the conversion. Three consequences follow. The nutrient is proportional to its held value, so the examples do not depend on the level of the supply. It depends on the composition only through the living volume fraction $\kappa$ of the field, not on which species is alive. And it is depleted over the length $\lambda=\ell/(\Lambda\sqrt\kappa)$, which is $1$~mm, the half edge of the cube, for fully living biofilm and $2.6$~mm when $\lv=0.15$: with $\Lambda=1$ the nutrient is never exhausted inside the biofilm, it decreases by a factor of order one across the cube. A slab of living biofilm with uniform $\kappa$ between the supply face and the opposite wall gives $c=\cosh((\ell-y)/\lambda)/\cosh(2\ell/\lambda)$ and at the centre $0.41$ for $\kappa=1$ and $0.82$ for $\kappa=0.15$; the seed means of Section~\ref{sec:composition} ($0.54$ and $0.91$) lie above these slab values because the biofilm outside the seed has $\phiD<1$ and does not fill the cube. The antibiotic field obeys the same equation with the same $\kappa$, so it is depleted over the same length; where living biofilm shields the nutrient it shields the antibiotic as well.

\paragraph{The antibiotic acts through a balance that depends on the density of the species.} The living fraction of species $i$ evolves by the second line of Eq.~\eqref{eq:point}, Eq.~17 of \citet{Klempt2026ContinuumBacterialGrowth}. At fixed $\phi_i$ and away from the barriers,
\begin{equation}
\eta_i\,\phi_i^2\,\dot\psi_i=c^*\phi_i\textstyle\sum_jA_{ij}\phi_j\psi_j-\alpha^*b_i\,\psi_i,
\label{eq:psi}
\end{equation}
so $\psi_i$ rises where the nutrient term exceeds the antibiotic term and falls where it does not. For a species that dominates its point ($A_{ii}\phi_i^2\psi_i\gg\sum_{j\neq i}A_{ij}\phi_j\psi_j$) the balance is $c^*A_{ii}\phi_i^2$ against $\alpha^*b_i$, and the relative rate is
\begin{equation}
\frac{\dot\psi_i}{\psi_i}=\frac{c^*A_{ii}}{\eta_i}-\frac{\alpha^*b_i}{\eta_i\,\phi_i^2}.
\label{eq:rate}
\end{equation}
The nutrient term gives a rate that does not depend on how much of the species is present, the antibiotic term a decay rate that grows as $1/\phi_i^2$: the antibiotic kills a species that fills its point slowly and the same species present with a small volume fraction almost at once. The rate changes sign on the line $c^*A_{ii}\phi_i^2=\alpha^*b_i$, but $\psi_i$ decays exponentially rather than reaching zero, and $\phi_i$ evolves at the same time by the first line of Eq.~\eqref{eq:point}, so the equations give a balance with a graded outcome, not a sharp boundary between living and dead regions. In the examples $c^*=c^*_0c$ and $\alpha^*=\alpha^*_0a$ with $c^*_0=\alpha^*_0$, so the balance for species~1 ($b_1=1$) reads $c\,\phi_1^2$ against $a$: it is approached only where the nutrient is as high as the antibiotic and the species fills its point, next to the supply face, and everywhere else the species declines at a rate set by $a/\phi_1^2-c$. Where the nutrient is high the decline is slow, where it is low it is fast; this is the spatial pattern of Section~\ref{sec:composition}. The same expression says that a mixed seed, in which each species holds about half the volume, loses its living fraction much faster than a seed of one species, and that the state given to newly reached material decides its fate under the antibiotic (Section~\ref{sec:newpoints}).

\paragraph{The time link is not a parameter of the model.} Both lines of Eq.~\eqref{eq:point} are linear in the rates, and every rate carries the viscosity of its species. Running the point model on the time $T^*$ of the field with $t=s\,T^*$ therefore equals running it on its own time with the viscosities $\eta_i/s$ and $\eta_{\phi,i}/s$ (the barrier penalty $K_{p,i}=10^{-4}\eta_i$ scales with them). With $T^*=1$ taken as the duration of the data the point model was calibrated on, $s$ is fixed by the data and the viscosities keep the values of the calibration; no new constant enters.

\subsection{Parameters}
\label{sec:parameters}

All constants are those of the papers (Table~\ref{tab:parameters}): Table~2 of \citet{Klempt2024DiffusionDrivenGrowth}, converted from their $20~\mu$m cube to the 2~mm cube of the examples, and the antibiotic of case~4 of \citet{Klempt2026ContinuumBacterialGrowth}, $\alpha^*_0=100$ from $T^*=1/3$. With $d=10^8$~mm$^2/T^*$ the fields adjust within about $4\times10^{-8}\,T^*$, so the fields are quasi-steady; keeping the time derivative of Eq.~\eqref{eq:35} changes the seed values by less than $0.005$. A diffusion slower by eight orders of magnitude ($d=1$~mm$^2/T^*$) is only a sensitivity: the antibiotic then reaches the seed more slowly and kills markedly less of it. Values not given by the papers are marked as assumptions: the length and time scale of the conversion, $\phi_{\min}$ (results change by less than $1\,\%$ below $0.001$), the start rule of new points, the void floor $f$ and the regularisation $\varepsilon$.

\begin{table}[tbp]
  \centering
  \caption{Constants of the examples. Conversion: Table~2 of \citet{Klempt2024DiffusionDrivenGrowth} on the 2~mm cube with the half edge $\ell=1$~mm, so that $g\ell^2/d=1$.}
  \label{tab:parameters}
  \footnotesize
  \setlength{\tabcolsep}{4pt}
  \begin{tabular}{lll>{\raggedright\arraybackslash}p{5.2cm}}
    \toprule
    quantity & symbol & value & source \\
    \midrule
    diffusion of $\phiD$ & $\beta$ & $0.02$~mm$^2/T^*$ & \citet{Klempt2024DiffusionDrivenGrowth}, converted (assumption) \\
    front speed, half saturation & $r$, $k$ & $10$~mm$/T^*$, $1$ & \citet{Klempt2024DiffusionDrivenGrowth}, converted \\
    growth rate & $k_\alpha$ & $10^{-3}\,T^{*-1}$ & \citet{Klempt2024DiffusionDrivenGrowth} \\
    nutrient diffusion, uptake & $d$, $g$ & $10^8$~mm$^2/T^*$, $10^8\,T^{*-1}$ & \citet{Klempt2024DiffusionDrivenGrowth}, converted \\
    antibiotic diffusion, uptake & $d_a$, $g_a$ & $=d$, $=g$ & assumption \\
    stiffness, Poisson ratio & $E$, $\nu$ & $10$~Pa, $0.49$ & \citet{Klempt2024DiffusionDrivenGrowth} \\
    void floor & $f$ & $10^{-3}$ & assumption \\
    point model & $\mathbf A$, $b$, $\eta$, $c^*_0$ & case~I & \citet{Fritsch2025BayesianMicrofilms} \\
    antibiotic level, onset & $\alpha^*_0$, $T^*_{\rm ab}$ & $100$, $1/3$ & \citet{Klempt2026ContinuumBacterialGrowth}, case~4 \\
    time link & $s$ & $0.1$ & from the data of case~I \\
    start threshold & $\phi_{\min}$ & $0.001$ & assumption (converged) \\
    front regularisation & $\varepsilon^2$ & $10^{-6}$ & assumption \\
    \bottomrule
  \end{tabular}
\end{table}

\section{Implementations and verification}
\label{sec:verification}

The coupled model is solved in three codes that share the equations, the parameters and the compiled point-model routine, but not the discretisation: a finite-volume code (7-point Laplacian, second-order upwind front term with the minmod limiter, RK4 in time, hexahedral mechanics with selective reduced integration), the neighboured-element user element of the original authors' group in ANSYS \citep{Blaszczyk2022TaylorWLS,Rudolf2025NEM}, in which the Laplacian and the gradients are fitted on the Gauss points from 30 neighbours by weighted least squares and the fields are updated explicitly, and a Galerkin user element in Abaqus (8-node hexahedra, SUPG \citep{BrooksHughes1982SUPG} for the front term, $\bar{\mathbf B}$ mechanics, implicit Euler with Newton's method). Three numerical details of the neighboured element were corrected before use, without changing its equations (Appendix~\ref{app:element}). Differences between the codes therefore come from the fields and the mechanics only.

The verification follows the usual steps \citep{Roache2002MMS,Celik2008GCI}; Appendix~\ref{app:verification} gives the details. With uniform fields every code reproduces the point model alone to round-off ($2\times10^{-14}$ in the finite-volume code, $1.8\times10^{-12}$ in the element), so the hand-over of $\lv$ to the consumption is exact. With a manufactured solution the finite-volume and the Galerkin code converge at second order in $L_2$ for $\phiD$ and $c$ (observed orders $1.7$ to $2.05$; the SUPG form of $\phiD$ at about $1.5$), and the analytic tangent of the Galerkin element agrees with central differences to $5\times10^{-11}$; the neighboured element does not converge next to the walls (Appendix~\ref{app:element}). The mechanics of two codes agree to $1\,\%$ on the same growth field, so differences in the stress between the codes come from the growth field.

One numerical property of the model sets the accuracy of the stress and is therefore stated here. The front term of Eq.~\eqref{eq:34} is an advection with the speed $r\,c/(k+c)$ along $\nc$ against the small diffusion $\beta$: the cell P\'eclet number $rh/(2\beta)$ is about $31$ on a $16^3$ grid of the 2~mm cube, so a central discretisation oscillates and needs stabilisation. A constant added diffusion $rh/4$, which keeps the element as delivered stable, smears the edge of the seed over a few cells; the seed stress, which is set by the jump of $\alpha$ across that edge, follows the smearing and is about a quarter of the converged value on $16^3$ and still rising on $96^3$. The growth variable converges with every treatment. The three second-order treatments (upwind with the minmod limiter, SUPG, upwind-biased least-squares gradients) give the same seed stress within a few per cent from $32^3$ on, and the three codes agree within their grid-convergence bands (Appendix~\ref{app:front}). All results below use a second-order treatment.

The examples were then run on $32^3$, $64^3$ and $128^3$ cells with the outer step $0.0025$ (the step study is in Appendix~\ref{app:timestep}) and the three-grid index of \citet{Celik2008GCI} with the safety factor $1.25$ (Figure~\ref{fig:grid}, Table~\ref{tab:gci}). The growth converges with a GCI of $0.7$ to $3.3\,\%$ in all examples, the seed stress with $0.7\,\%$ without and $8.3\,\%$ with the antibiotic when the species start apart. When they start mixed the stress converges slowly (observed order about $0.75$, GCI $26$ to $29\,\%$), because the front carries intermediate values of $\phiD$ at a nearly discontinuous edge and the largest $\phiD$ in the domain still rises with the grid ($0.60$, $0.69$, $0.73$); with the start fractions of \citet{Fritsch2025BayesianMicrofilms} the $128^3$ grid is not yet in the asymptotic range. No run on $32^3$ to $128^3$ left $[0,1]$ by more than $2\times10^{-4}$.
\TODO{Fourth grid ($192^3$, step limit of Appendix~\ref{app:timestep}) for the mixed example with the start fractions of Fritsch et al.; add to Table~\ref{tab:gci} and Figure~\ref{fig:grid} when available.}

\begin{figure}[tbp]
  \centering
  \includegraphics[width=0.75\linewidth]{\figs fig_grid_examples.pdf}
  \caption{Grid convergence of the two-species examples (finite-volume code, outer step $0.0025$; Table~\ref{tab:gci}): seed means of (a)~$\alpha-1$ and (b)~the von Mises stress at $T^*=1$ against the grid. Dashed: Richardson extrapolation where the three grids are in the asymptotic range.}
  \label{fig:grid}
\end{figure}

\section{Results}
\label{sec:results}

\subsection{Set-up}
\label{sec:setup}

The examples use a 2~mm cube with a cubic seed of $\phiD=1$ at its centre, the nutrient held at $c=1$ on the face $y=-1$~mm and the antibiotic held at $a=1$ on all outer faces from $T^*=1/3$. In the separate example species~A fills the seed and species~B four small blocks along the edges of the cube; in the mixed example both start in the seed, in equal amounts ($\phi_1=\phi_2=0.275$) or in the proportions of \citet{Fritsch2025BayesianMicrofilms} ($0.25$, $0.30$). The front term is on in every run. The values below are seed means at $T^*=1$ from the Python code (TVD, newly reached points inherit the living composition of their neighbours), which reaches the finest grids; the element gives the same composition within its grid error and its growth within $4\,\%$ after extrapolation. The stress is solved on $16^3$ hexahedra with the element's integration; solving the same $64^3$ field on $32^3$ and $64^3$ elements changes the seed von Mises stress by $0.2\,\%$.

\subsection{Composition and living biomass}
\label{sec:composition}

Table~\ref{tab:examples} gives the seed means and Figure~\ref{fig:panels} the fields of the separate and the mixed example with antibiotic. Without antibiotic the living volume fraction of the seed stays near $0.96$; the antibiotic lowers it to $0.15$ when the species start apart and to $0.04$ when they start mixed. Species~B is removed first: its nutrient term $c^*A_{22}/\eta_2$ is half that of A, while the antibiotic term per viscosity, $\alpha^*b_i/\eta_i$, is the same for both. The nutrient field decides where the antibiotic acts, as Eq.~\eqref{eq:rate} says: species~A declines at the rate $(a/\phi_1^2-c)\,c^*_0s/\eta_1$ per unit of $T^*$ ($\eta_1=1$), so where the nutrient is high it survives, and where it is low it dies as well. The seed values give the size of the rates: in the separate example the seed holds $\phi_1=0.91$ with $c=0.91$ and $a=0.98$ at $T^*=1$, a decline of $2.7$ per unit of $T^*$, consistent with the fall of $\psi_1$ from $0.98$ to $0.16$ between the onset and $T^*=1$; in the mixed seed species~A holds only $\phi_1=0.57$, the decline is $21$ per unit of $T^*$, and its living fraction is gone within $0.2\,T^*$ of the onset. The difference between $0.15$ and $0.04$ is therefore a property of the point model's balance, not of the field: the same species at the same nutrient and antibiotic dies ten times faster when it shares its point. When the biomass dies the nutrient is no longer consumed and rises in the seed, from $0.54$ to $0.91$ in the separate example, as the depletion length of Eq.~\eqref{eq:thiele} grows from $1$ to $2.6$~mm; the antibiotic, which obeys the same equation, reaches $0.98$ in the seed. Living volume fraction, shares and nutrient change by less than $3\,\%$ between the grids ($32^3$ to $128^3$) and the outer steps ($0.02$ to $0.0025$) used.

\begin{table}[tbp]
  \centering
  \caption{Seed means at $T^*=1$, Python, $64^3$ cells, outer step $0.025$; converged growth and stress in Table~\ref{tab:gci}.}
  \label{tab:examples}
  \small
  \setlength{\tabcolsep}{4pt}
  \begin{tabular}{lccccc}
    \toprule
    example & $\lv$ & $\phi_1/(\phi_1+\phi_2)$ & $c$ & $\alpha-1$ & von Mises [Pa] \\
    \midrule
    separate & $0.965$ & $0.978$ & $0.536$ & $9.06\times10^{-4}$ & $6.56\times10^{-4}$ \\
    separate, antibiotic & $0.152$ & $0.951$ & $0.911$ & $8.95\times10^{-4}$ & $5.68\times10^{-4}$ \\
    mixed, equal amounts & $0.961$ & $0.859$ & $0.608$ & $8.33\times10^{-4}$ & $3.50\times10^{-4}$ \\
    mixed, equal amounts, antibiotic & $0.038$ & $0.610$ & $0.973$ & $8.30\times10^{-4}$ & $3.23\times10^{-4}$ \\
    mixed $(0.25,0.30)$, antibiotic & $0.037$ & $0.510$ & $0.973$ & $8.30\times10^{-4}$ & $3.22\times10^{-4}$ \\
    \bottomrule
  \end{tabular}
\end{table}

\begin{figure}[p]
  \centering
  \includegraphics[width=0.96\linewidth]{\assets fig_coupled_3d_panels_separate_ab.pdf}\\[1mm]
  \includegraphics[width=0.96\linewidth]{\assets fig_coupled_3d_panels_mixed_ab.pdf}
  \caption{Separate (upper half) and mixed (lower half) example with antibiotic, Python, $64^3$, one octant removed, at $T^*=0.325$ (first row of each half) and $T^*=1$ (second row): (a)~share $\phi_1/(\phi_1+\phi_2)$, (b)~living volume fraction $\lv$, (c)~nutrient $c$, (d)~antibiotic $a$, (e)~von Mises stress. In the mixed example the living volume fraction falls almost everywhere after the antibiotic.}
  \label{fig:panels}
\end{figure}

\subsection{Antibiotic over time}
\label{sec:schedule}

With the same total dose given constantly, in pulses or as a sine from $T^*=1/3$ (separate example, $64^3$; the time series are in the supplementary material), $\alpha-1$ changes by $0.1\,\%$ and the seed stress by $0.8\,\%$ between the three schedules, while the living volume fraction of the seed follows the dose: in the pauses of the pulsed dose it recovers from $0.17$ to up to $0.34$ and ends at $0.30$, and the sine gives a lower minimum ($0.11$ against $0.15$) and a higher end value ($0.19$) than the constant dose ($0.15$). How the dose is given acts on the composition and the living biomass, hardly on the mechanics.



\subsection{Which composition does new biofilm receive?}
\label{sec:newpoints}

The rule for newly reached points (Section~\ref{sec:coupling}, item~4) is not fixed by the papers: the point model has no transport term, Eq.~\eqref{eq:34} moves $\phiD$ without carrying the point state, so the composition travels with the biofilm only through the state that a newly reached point is given. The two rules compared here differ in the volume fraction they assign. Inheriting the neighbours' amounts, shared by their living biomass, gives a new point a volume fraction close to that of the established biofilm behind it; the start state of the colony gives it the small start fractions of the colony ($0.25$ to $0.3$ per species). By Eq.~\eqref{eq:rate} the antibiotic term scales as $1/\phi_i^2$, so a point started from the start state loses its living fraction at a rate about ten times that of an inherited point ($a/\phi_1^2$ of $16$ against about $1.2$). Where the biofilm has stopped spreading the rule hardly matters: with the start values of the seed instead of the living composition of the neighbours, the seed values of the cube examples change by a few per cent ($\lv$ $0.147$ against $0.152$ in the separate example with antibiotic). In the maze of \citet[Fig.~11]{Klempt2024DiffusionDrivenGrowth}, where the biofilm keeps spreading after the antibiotic and most of the colonised area at the end was reached after the onset, the rule decides the outcome (Figure~\ref{fig:mazerules}, set-up in Section~\ref{sec:geometry}): the living volume fraction of the colonised region falls to $0.10$ shortly after the onset of the antibiotic at $T^*=20$ and recovers to $0.35$ at $T^*=60$ when new cells take the state of their neighbours, because the biofilm keeps spreading with dense living cells that the antibiotic removes slowly, and it stays at $0.04$ when they start from the start state of the colony, because thinly started cells die at once. The Galerkin element with the start-state rule gives the same (0.04 to 0.07 in both codes). Without antibiotic the two codes agree in the colonised area within $3$ to $5\,\%$ and in the living fraction within $0.006$.

\begin{figure}[tbp]
  \centering
  \figorbox{\figs fig_maze_rules.pdf}{width=0.75\linewidth}
  \caption{The maze with the antibiotic held at the entrance from $T^*=20$, at $T^*=60$: living volume fraction $\lv$ where $\phiD\ge0.05$ when newly reached cells (a)~take the state of their colonised neighbours and (b)~take the start state of the colony; walls grey, nutrient source outlined in blue, antibiotic source in red. (c)~Living volume fraction averaged over the colonised cells against $T^*$ for both rules; the antibiotic starts at the dotted line.}
  \label{fig:mazerules}
\end{figure}

\subsection{What reaches the stress}
\label{sec:stress}

The growth $\alpha-1$ hardly depends on the composition, as the bound of Section~\ref{sec:structure} requires: Eq.~\eqref{eq:36} is unchanged, the antibiotic lowers $\psi_i$ but not $\phiD$, and the seed growth is within $10\,\%$ of the bound $k_\alpha T^*$ when the species start apart and within $15\,\%$ when they start mixed (Table~\ref{tab:gci}: $9.16$ against $9.22\times10^{-4}$ in the separate example, about $1\,\%$ between the two cases; $8.56$ against $8.54\times10^{-4}$ in the mixed one). The composition reaches the stress along the path~\eqref{eq:path} only, through the front speed, which the nutrient rise from $0.54$ to $0.91$ raises by $37\,\%$ at the edge of the seed, but in these examples the effect of the antibiotic on the converged seed stress is within the numerical uncertainty: $6.55$ without and $6.70\times10^{-4}$~Pa with the antibiotic in the separate example (extrapolated, GCI $0.7$ and $8.3\,\%$). On a fixed grid the difference is larger ($6.73$ against $6.00\times10^{-4}$~Pa on $64^3$) because the two cases converge from opposite sides. The mixed seed carries about half the stress of the separate one, because its edge is less sharp: the front carries intermediate values of $\phiD$. A limitation follows from the same structure: the front speed $r\,c/(k+c)$ depends on the nutrient only, so dead biofilm keeps spreading after the antibiotic; in the maze about two thirds of the cells reached after the antibiotic hold only dead biofilm when new cells start from the start state. Tying the front speed or the growth rate to the living biomass would open a second path and changes the model (Section~\ref{sec:discussion}).

\subsection{A slower front}
\label{sec:frontspeed}

With the front speed $r=10$~mm$/T^*$ of \citet{Klempt2024DiffusionDrivenGrowth}, converted to the 2~mm cube, the biofilm reaches the walls before the antibiotic starts and fills the cube by $T^*=1$: $92$ to $95\,\%$ of the cells hold $\phiD\ge0.5$ at the end. To check that the results above are not a property of a full cube, the four examples were repeated with one tenth of that speed, $r=1$~mm$/T^*$, all else unchanged ($64^3$, outer step $0.025$; Table~\ref{tab:slowfront}). The biofilm then occupies $20$ to $27\,\%$ of the cube at $T^*=1$ and is still spreading. The seed means change in the way the path of Eq.~\eqref{eq:path} prescribes. The nutrient in the seed is higher without antibiotic ($0.75$ against $0.54$ and $0.61$), because less living biofilm lies between the supply face and the seed, so the antibiotic balance of Eq.~\eqref{eq:rate} tips less far in the separate seed, which keeps a living volume fraction of $0.30$ instead of $0.15$, while the mixed seed, whose species share their points, loses its living biomass as before ($0.04$); the higher nutrient also lets species~A take over the mixed seed further ($0.98$ against $0.86$ without antibiotic). The growth $\alpha-1$ of the seed stays within $8\,\%$ of its value at the full speed, as the bound $k_\alpha T^*$ requires, and the antibiotic changes it by $3\,\%$ (separate) and $0.2\,\%$ (mixed) at the slower speed, against about $1\,\%$ at the full speed. The seed stress is $1.6$ to $3$ times higher at the slower front, because the edge of the seed stays sharper when it moves less far, and on this grid the antibiotic changes it by $15\,\%$ (separate, upwards) and $6\,\%$ (mixed), the same order as the $16$ and $8\,\%$ on the same grid at the full speed, where the extrapolation of Table~\ref{tab:gci} reduces the difference to the grid uncertainty; a grid series at the slower speed was not run. The front speed sets how much of the domain the biofilm occupies, how high the stress of the seed is and how far the antibiotic tips the balance in the seed; what the composition does to the growth does not depend on it.

\begin{table}[tbp]
  \centering
  \caption{One tenth of the front speed ($r=1$~mm$/T^*$) against the speed of the paper ($r=10$~mm$/T^*$): seed means at $T^*=1$ and the colonised fraction of the cube ($\phiD\ge0.5$), Python, $64^3$ cells, outer step $0.025$; mixed: equal amounts; $r$ in mm$/T^*$.}
  \label{tab:slowfront}
  \footnotesize
  \setlength{\tabcolsep}{3pt}
  \begin{tabular}{llcccccc}
    \toprule
    example & $r$ & $\lv$ & $\phi_1/(\phi_1+\phi_2)$ & $c$ & $\alpha-1$ & von Mises [Pa] & colonised \\
    \midrule
    separate & $10$ & $0.965$ & $0.978$ & $0.536$ & $9.06\times10^{-4}$ & $6.56\times10^{-4}$ & $0.95$ \\
    separate & $1$ & $0.969$ & $0.981$ & $0.745$ & $8.46\times10^{-4}$ & $1.04\times10^{-3}$ & $0.23$ \\
    separate, antibiotic & $10$ & $0.152$ & $0.951$ & $0.911$ & $8.95\times10^{-4}$ & $5.68\times10^{-4}$ & $0.95$ \\
    separate, antibiotic & $1$ & $0.299$ & $0.965$ & $0.918$ & $8.70\times10^{-4}$ & $1.19\times10^{-3}$ & $0.27$ \\
    mixed & $10$ & $0.961$ & $0.859$ & $0.608$ & $8.33\times10^{-4}$ & $3.50\times10^{-4}$ & $0.92$ \\
    mixed & $1$ & $0.969$ & $0.981$ & $0.756$ & $9.03\times10^{-4}$ & $1.14\times10^{-3}$ & $0.22$ \\
    mixed, antibiotic & $10$ & $0.038$ & $0.610$ & $0.973$ & $8.30\times10^{-4}$ & $3.23\times10^{-4}$ & $0.92$ \\
    mixed, antibiotic & $1$ & $0.039$ & $0.658$ & $0.984$ & $9.01\times10^{-4}$ & $1.08\times10^{-3}$ & $0.20$ \\
    \bottomrule
  \end{tabular}
\end{table}

\subsection{Geometry: rigid bars and a maze}
\label{sec:geometry}

\citet{Klempt2024DiffusionDrivenGrowth} also grow the biofilm against rigid obstacles (their Section~4.3) and through a maze (Section~4.4). Both were repeated with the coupled model in the Python code, without antibiotic unless stated (Figure~\ref{fig:geometry}).

\paragraph{Rigid bars.} Four bars after Fig.~9 of the paper, with positions read off the sketch: two along $x$ and two along $z$, each $0.25$~mm square, between $0.5$ and $0.75$~mm from the centre. They hold no biofilm and no nutrient, and the biofilm slides along them without friction. All other settings are those of Section~\ref{sec:setup} ($32^3$, outer step $0.0025$). The biofilm grows around the bars and closes behind them, and by $T^*=0.2$ it fills the cube between them. The composition of the seed hardly changes. The bars restrain the swelling: the seed von Mises stress is $7.1\times10^{-3}$~Pa against $7.4\times10^{-4}$~Pa without bars in the separate example, about ten times higher, and $9.5\times10^{-4}$ against $2.1\times10^{-4}$~Pa in the mixed one; the growth $\alpha-1$ of the seed is $9$ and $21\,\%$ lower. These are values on one grid; on $32^3$ the stress without bars lies $14\,\%$ above its extrapolated value (Table~\ref{tab:gci}), small against the factor caused by the bars.

\paragraph{Maze.} The maze of Fig.~11 of the paper, read off the figure ($560~\mu$m across), in two dimensions with $4~\mu$m cells, the constants of Table~6 of the paper ($\beta=10~\mu$m$^2/T^*$, $r=200~\mu$m$/T^*$, $k_\alpha=10^{-6}\,T^{*-1}$), the nutrient held in the centre, the colony of both species at the entrance and the point model of the two-species examples. With the antibiotic, $a=1$ is held at the entrance from $T^*=20$. The model is two-dimensional, so it gives no stress. The biofilm follows the corridors, also into dead ends, and approaches the centre: at $T^*=60$ the nearest point with $\phiD\ge0.5$ is $42~\mu$m from the centre without antibiotic and $133~\mu$m with it. The living volume fraction of the colonised region is $0.92$ without antibiotic; with it, the rule for new cells decides whether the biofilm recovers (Section~\ref{sec:newpoints}).

\begin{figure}[tbp]
  \centering
  \includegraphics[width=0.55\linewidth]{\assets fig_obstacle_columns_3d.pdf}\\[1mm]
  \includegraphics[width=0.7\linewidth]{\assets fig_maze_k24_h4.pdf}
  \caption{Upper half: growth against four rigid bars (dark grey) after Fig.~9 of \citet{Klempt2024DiffusionDrivenGrowth}, Python, $32^3$, no antibiotic. (a)~Set-up: start biomass (red) and nutrient supply on the face $y=-1$~mm (green). (b)--(d)~Surface $\phiD=0.5$ (separate example) and $0.2$ (mixed example, whose front carries lower values of $\phiD$) at $T^*=0.05$, $0.1$ and $0.2$. Lower half: the maze of Fig.~11 of the same paper with the coupled model (two dimensions, $4~\mu$m cells): $\phiD$ at $T^*=12$, $24$, $36$ and $60$, (a)~without and (b)~with the antibiotic at the entrance from $T^*=20$. Walls white, nutrient source outlined in red.}
  \label{fig:geometry}
\end{figure}

\subsection{A clinical molar crown}
\label{sec:realtooth}

The cube examples place the seed in a box. To show the coupling on a real geometry, the same model and constants were run around the crown of a human molar from the Open-Full-Jaw dataset \citep{Gholamalizadeh2022OpenFullJaw} (Figure~\ref{fig:realtooth}). The crown ($6.7\times6.8\times5.5$~mm) stands with its cervical base on the bottom face of a 10~mm box and is rigid for the fields (no biofilm and no nutrient inside, zero flux through its surface) and for the stress ($\mathbf u=\mathbf 0$). Nutrient and, from $T^*=1/3$, antibiotic are held at $1$ on the five faces around the crown (the saliva side); the bottom face is zero flux (the root side). The biofilm starts as a cervical ring $0.25$~mm thick and $1$~mm high, species~A on the half $x<0$ and B on $x>0$. These choices are assumptions of this work. The grid has $80^3$ cells ($h=0.125$~mm; $40^3$ for comparison), the outer step is $0.025$, the film is taken as the cells with $\phiD>0.5$, and the stress is solved on the field grid with the film surface free.

With the paper's front speed the biofilm spreads along the bottom face and outwards and reaches the side faces by $T^*=1$. The film that covers the bottom face around the crown at the end is a consequence of that speed, which fills the 2~mm cube in the same time (Section~\ref{sec:frontspeed}), not of the composition. The composition is carried with it: the A and B halves stay separate (film mean share $0.50$; $0.956$ in the A ring, $0.051$ in the B ring). The antibiotic enters with the nutrient ($a=c$ in the film, $0.73$ on average, $0.24$ at the ring) and acts first on B, as $b_2=2$ and Eq.~\eqref{eq:rate} require: the living fraction of the B ring falls from $0.71$ at the onset to $0.13$ at $T^*=1$, that of the A ring from $0.81$ to $0.57$; the film mean is $0.37$ ($0.95$ without antibiotic). Growth and stress hardly depend on the antibiotic, as on the cube (film mean $\alpha-1$ of $4.3\times10^{-4}$ with and without; von Mises $3.0$ and $3.2$~mPa, with the largest values in the corners between the film, the crown and the bottom face). Between $40^3$ and $80^3$ the film means change by $3\,\%$ and the film stress by $7\,\%$, most of which is the stress mesh. The example demonstrates the coupling on a real surface; the stress of the cervical region with a growth law is not treated here.

\begin{figure}[tbp]
  \centering
  \includegraphics[width=0.8\linewidth]{\figs fig_real_tooth_coupled.pdf}
  \caption{The coupled model around the crown of a human molar \citep{Gholamalizadeh2022OpenFullJaw}, $80^3$ cells, with the antibiotic from $T^*=1/3$. (a)~Start: the cervical ring, species~A on the half $x<0$ (blue), B on $x>0$ (red). (b)--(f)~The film $\phiD>0.5$ at $T^*=1$, coloured by (b)~the share $\phi_1/(\phi_1+\phi_2)$, (c)~the living volume fraction $\lv$, (d)~the nutrient $c$, (e)~the antibiotic $a$ and (f)~the von Mises stress.}
  \label{fig:realtooth}
\end{figure}

\section{Discussion}
\label{sec:discussion}

\paragraph{What the coupling shows.} The examples confirm what the structure of the equations prescribes (Section~\ref{sec:structure}). The composition reaches the growth field only through the consumption $g\,\phiD(\lv)\,c$, and from there only through the front speed $r\,c/(k+c)$. When the antibiotic kills the biomass, less nutrient is consumed, the nutrient rises in the seed (its depletion length grows from $1$ to $2.6$~mm) and the edge moves faster, by $37\,\%$ at the nutrient levels of the seed. The effect on the growth is about $1\,\%$ and could not exceed about $10\,\%$, because the growth law $\dot\alpha=k_\alpha\phiD$ counts all biomass, living or dead, and the seed is already near the bound $k_\alpha T^*$; the effect on the seed stress is within its grid uncertainty, because the stress is set by the jump of $\alpha$ across the edge of the seed and the edge moves but does not change its height. The local nutrient decides where the antibiotic acts, and the balance of Eq.~\eqref{eq:rate} says how: a species declines at a rate $\alpha^*b_i/(\eta_i\phi_i^2)-c^*A_{ii}/\eta_i$, slow where the nutrient is high and the species is dense, fast where the nutrient is low or the species is thinly present. The same dose therefore removes a species in one part of the biofilm and not in another, which a model without space cannot show, and removes it ten times faster from a mixed seed than from a seed of one species, which a model without the point model's balance cannot show either. Numerically, the stress depends on how sharp the edge of the biofilm is, and that edge is set by the front term, so the coupling needs a second-order treatment of that term (Section~\ref{sec:verification}).

\paragraph{What a composition-dependent mechanics needs.} The single path of Eq.~\eqref{eq:path} is the main result of this work, and its consequence is a statement about models rather than about these examples: with the published equations the composition cannot change the stress by more than the front speed allows, whatever the species and whatever the antibiotic does to them, so a model in which the composition is to change the mechanics needs a second path. The coupling presented here is the base on which such a path can be added without touching the hand-over of $c^*$, $\alpha^*$ and $\lv$, and the runs give an estimate of what each would do. (i)~Growth from the living biomass, $\dot\alpha=k_\alpha\phiD\lv$ as in \citet{WannerGujer1986}, would make the antibiotic act on the growth directly. Its size can be estimated from the time series of the seed without a new run, since $\alpha-1$ would then scale with the time average of $\lv$ over $[0,T^*]$: $0.54$ against $0.92$ with and without the antibiotic in the separate example and $0.30$ against $0.88$ in the mixed one, so the seed growth would fall by about $40$ and $65\,\%$ in place of the $1\,\%$ of the published law. This is an estimate from the published runs, not a computed value. (ii)~A front speed tied to the living biomass, which would also stop dead biofilm from spreading after the antibiotic, changes the stress through the sharpness of the edge: in a test on $16^3$ (front speed multiplied by the local living fraction) the seed stress rose by $51\,\%$ in the separate example and about fourfold in the mixed one. (iii)~A stiffness that depends on the species changes the seed stress by $9$ to $89\,\%$ in a $16^3$ test, depending on which species is stiffer. Each of these is a change of the model, and the value of the coupling as presented is that it separates what the published equations give from what such a change would add.

\paragraph{Comparison with experimental scales.} The examples are verified against each other and against the equations, not validated against measurements, and the scales of the model can be set against measured ones only loosely. Thickness: on enamel slabs worn in situ, dental plaque reaches $14.9\pm5.0~\mu$m after one day and $49.3\pm11.6~\mu$m after seven days \citep{AlAhmad2007PlaqueFISH}. The $20~\mu$m cube of \citet{Klempt2024DiffusionDrivenGrowth} is of the order of a one-day plaque; the 2~mm cube of the examples, chosen for the element, and the $0.25$~mm ring around the molar crown are 40 and 5 times the seven-day thickness. Stiffness: the compressive elastic modulus of in vitro \emph{S.~mutans} biofilms is $0.28$~kPa (standard deviation $0.35$~kPa, $n=30$) \citep{Rmaile2013MicrobialTribology}, about 30 times the $E=10$~Pa of \citet{Klempt2024DiffusionDrivenGrowth}; since the mechanics is coupled one way, every stress of this work scales linearly with $E$, so the seed stress of $6.6\times10^{-4}$~Pa becomes about $0.02$~Pa at that modulus and the stress between rigid bars about $0.2$~Pa. Stress at which growth stops: confined colonies of \emph{E.~coli} arrest their growth at pressures of about $10$~kPa or below \citep{Chu2018SelfInducedStress}, and a solid stress of $6$ to $16$~kPa inhibits the growth of tumour spheroids \citep{Helmlinger1997SolidStress}. The growth stresses of the examples lie four to seven orders of magnitude below these values, which is consistent with neglecting the mechanical feedback on the growth in this regime, but it is not a test of the model: the growth of the examples is $\alpha-1\approx10^{-3}$, and a plaque that doubles its thickness carries a growth of order one. The time unit $T^*$ is fixed by neither paper, and the point model's data are synthetic. A validation would need, for one oral multi-species biofilm, the thickness over days (which fixes $T^*$ and $k_\alpha$), a measured modulus of the intact film, and a measured deformation or stress of the growing film on its substrate; none of these was found for oral plaque in the sources checked.

\paragraph{Limitations.} (i)~The front speed depends on the nutrient only, so dead biofilm keeps spreading after the antibiotic, and with the front speed of \citet{Klempt2024DiffusionDrivenGrowth} the biofilm fills the cube within $T^*=1$; at one tenth of that speed (Section~\ref{sec:frontspeed}) the seed stress is $1.6$ to $3$ times higher and the antibiotic tips the balance of the separate seed less far, while its effect on the growth stays within $3\,\%$. (ii)~Eq.~\eqref{eq:36} lets dead biomass grow at the same rate as living biomass. (iii)~The point model has no space; the composition moves only through the rule for newly reached points, which decides in the maze whether the biofilm recovers. A species-resolved continuum would replace the rule. (iv)~The mechanics is coupled one way, as in \citet{Klempt2024DiffusionDrivenGrowth}; with $E=10$~Pa the neglected term acts only for coefficients far beyond the paper values (an effect on the growth from about $10^5$ on, $16^3$ test). (v)~Living and dead biomass, and both species, have the same stiffness. (vi)~The examples use synthetic data with known parameters, the stress converges slowly when the species start mixed, and the element's values carry a grid error of $16.5\,\%$. Calibrated data \citep[for example the five-species data of][]{Heine2025PeriImplant} and a measured stress are needed before the stresses can be validated, not only verified.

\paragraph{Outlook.} Species-resolved transport, so that the composition moves with the biomass; front speed and growth tied to the living biomass, with the consequences for the stress quantified above; growth laws with a mechanical feedback; calibrated multi-species data in the element; and a realistic tooth or implant geometry. Each can be added to the coupling as presented here, and the verification procedure, several codes, a manufactured solution and grid convergence indices, carries over.

\section{Conclusions}
\label{sec:conclusions}

This work coupled the multi-species point model of \citet{Klempt2026ContinuumBacterialGrowth} to the continuum growth model of \citet{Klempt2024DiffusionDrivenGrowth} at every integration point, without changing the equations of either model. The point model receives the local nutrient and antibiotic from the fields, $c^*=c^*_0c$ and $\alpha^*=\alpha^*_0a$, and its living volume fraction sets the uptake $g\,\phiD(\lv)\,c$ of both.
\begin{enumerate}
\item In the coupled equations the composition reaches $\phiD$, $\alpha$ and the stress along one path only, through the nutrient and the front speed $r\,c/(k+c)$, because the growth law $\dot\alpha=k_\alpha\phiD$ counts all biomass; the seed growth is bounded by $k_\alpha T^*$ and lies within $10\,\%$ of it, and the front speed changes by at most a factor of two. The examples confirm it: the antibiotic changes the growth by about $1\,\%$ and the converged seed stress within its grid uncertainty ($6.55$ against $6.70\times10^{-4}$~Pa extrapolated in the separate example), in the cube and around a molar crown; at one tenth of the front speed of the paper the antibiotic changes the growth by at most $3\,\%$. The consequence is that any model in which the composition is to change the mechanics needs a second path: growth or front speed tied to the living biomass, or a species-dependent stiffness. The coupling presented here is the base for such a term, and the runs estimate its size: growth from the living biomass would lower the seed growth by about $40$ to $65\,\%$ (an estimate from the time series), a front speed tied to the living biomass raises the seed stress by $51\,\%$ to fourfold, a species-dependent stiffness by $9$ to $89\,\%$ ($16^3$ tests).
\item The quasi-steady nutrient depends on the composition only through $\phiD\lv$ and, with the Thiele modulus $\Lambda^2=g\ell^2/d=1$, is depleted over the half edge of the domain but never exhausted. The time link $s$ is equivalent to the viscosities $\eta_i/s$ and is fixed by the data.
\item In the point model the antibiotic acts through a balance, $c^*A_{ii}\phi_i^2$ against $\alpha^*b_i$, whose decay rate grows as $1/\phi_i^2$: a species that fills its point is removed slowly, the same species thinly present at once. This explains three results of the examples: the local nutrient decides where the antibiotic acts; a mixed seed loses its living biomass ten times faster than a seed of one species ($0.96$ to $0.04$ against $0.96$ to $0.15$); and the state given to newly reached material, which neither paper fixes, decides in a maze whether the biofilm recovers after the antibiotic (living volume fraction $0.35$ against $0.04$).
\item The coupling was implemented in a finite-volume code, in a neighboured-element user element in ANSYS and in a Galerkin user element in Abaqus, and verified by replay of the point model, a manufactured solution and grid convergence indices. The front term \fterm{} sets the accuracy of the stress: a constant added diffusion gives a quarter of the converged seed stress, a second-order treatment is needed; with it the growth converges with a GCI of $0.7$ to $3.3\,\%$ and the seed stress with $0.7$ to $8.3\,\%$ when the species start apart.
\end{enumerate}
The next steps are models in which the front speed, the growth or the stiffness depend on the living biomass or on the species, with the consequences estimated here, and a measured thickness, modulus and deformation of one oral biofilm, without which the stresses remain verified but not validated.

\section*{Data availability}
The finite-volume code, the point-model routine, the run scripts, the result files of the examples and the figure scripts are in the public repository \url{https://github.com/keisuke58/pde-fem-biofilm}. The examples use the published synthetic data of \citet{Fritsch2025BayesianMicrofilms}; no calibration data are used in this work. The molar crown of Section~\ref{sec:realtooth} is derived from the Open-Full-Jaw dataset of \citet{Gholamalizadeh2022OpenFullJaw}, licensed CC BY-NC-SA, and the derived surface is redistributed in the repository under the same licence. The user elements depend on the commercial codes and on the element of \citet{Rudolf2025NEM}; the material routines are in the repository.
\TODO{Confirm with all authors which parts of the element code may be published, and whether a release (DOI) of the repository is made at submission.}

\section*{Declarations}
Author list, contributions, funding and conflicts of interest: to be agreed.
\TODO{Authors, affiliations, contributions, funding (LUH/IKM, Keio), acknowledgements; apply the Springer Nature template (sn-jnl.cls) of Biomechanics and Modeling in Mechanobiology at submission.}

\appendix

\section{Numerical details}
\label{app:numerics}

\subsection{The neighboured element}
\label{app:element}

Three details of the element were corrected, without changing its equations: the surface points take the value of the nearest Gauss point instead of a distance-weighted mean, the Laplacian is used in a conservative form, and the neighbours of the fit are chosen by distance instead of by element number. Before, $\phiD$ leaked through the walls of the cube (the surface points took an inverse-square mean, and the column sums of the Laplacian were not zero), and the stencils of the $16^3$ and $24^3$ grids were asymmetric inside the domain. With a manufactured solution the corrected element still does not converge from $8^3$ to $16^3$ ($L_2$ error of $\phiD$ $2.2\times10^{-3}$ to $1.8\times10^{-3}$, of $c$ $8.2\times10^{-3}$ to $1.2\times10^{-2}$): the MLS stencils of the Gauss points next to a face include the surface points and are ill-conditioned (error of the Laplacian of the manufactured $c$: $0.88$ on $8^3$, $65$ on $16^3$, against $|\nabla^2c|=3.5$). Rebuilding these stencils with the neighbours mirrored across the faces (zero flux built in) makes the operators consistent up to the walls, and with the weight width proportional to $h$ the element converges ($\phiD$ $9.5\times10^{-4}$ to $3.4\times10^{-4}$, order $1.5$; $c$ $9.0\times10^{-4}$ to $2.2\times10^{-4}$, order $2.0$). In the coupled example the mirrored operators run on $8^3$ but become unstable at a corner of the $16^3$ grid after the antibiotic onset; they are not used for the results. The explicit update of the element is limited by $\beta\Delta t/h^2\le0.12$ to $0.16$ on $16^3$ and $24^3$ (Bloch analysis of the stencil); the Lax--Wendroff form with the extremum sensor keeps $\phiD$ in $[-0.03,1.01]$ on $64^3$ in Python but is unstable at the corners of the supply face on $24^3$ in the element, so the element's stresses are quoted from the upwind MLS gradients. The time-step series of the element and the stability analysis will be given in the supplementary material.

\section{Verification details}
\label{app:verification}

\subsection{Manufactured solution and mechanics}
\label{app:mms}

Sources for $\phiD$ and $c$ were built from a manufactured solution \citep{Roache2002MMS} by automatic differentiation and checked symbolically (agreement $10^{-14}$). The finite-volume and the Galerkin code converge at second order in $L_2$ (the error plot is in the supplementary material): observed orders between $8^3$, $16^3$ and $32^3$ of $1.85$ and $1.90$ for $\phiD$ and $2.03$ and $2.01$ for $c$ in Python, $1.68$ and $1.81$ for $\phiD$ and $2.00$ and $2.05$ for $c$ in Abaqus; on $32^3$ the errors of $\phiD$ are $2.0\times10^{-4}$ and $2.9\times10^{-4}$. For the SUPG form the order of $\phiD$ is about $1.5$: the largest error sits in a thin layer next to the wall opposite the supply, where $\nabla c\to\mathbf 0$ and the front direction changes over a few cells. The analytic tangent of the Abaqus element agrees with central differences to $5\times10^{-11}$. The last growth field of an Abaqus run, solved by the Python mechanics, gives the seed von Mises stress to $1.1\,\%$ (no obstacles) and $0.5\,\%$ (obstacles with frictionless contact).

\subsection{The front term and the agreement of the codes}
\label{app:front}

Four treatments of the front term were compared, none of which changes Eqs.~\eqref{eq:34} to \eqref{eq:36}: (1)~a constant added diffusion $r h/4$ (first order; the element as delivered needs it to stay stable); (2)~TVD upwinding (MUSCL with the minmod limiter; Python); (3)~SUPG with $\tau=\tfrac{h}{2|v|}(\coth Pe-1/Pe)$ on the strong residual (Abaqus); (4)~for the explicit element, either the Lax--Wendroff term $\tfrac{\Delta t}{2}\,v\cdot\nabla(v\cdot\nabla\phiD)$ with an added diffusion switched on by an extremum sensor, or upwind-biased MLS gradients. Figure~\ref{fig:convergence} shows the seed means of $\alpha-1$ and of the von Mises stress at $T^*=1$ of the separate example with antibiotic against the grid for every scheme and code. The growth variable converges for every scheme, and the codes agree within their grid-convergence bands: Python extrapolates to $9.0$ to $9.8\times10^{-4}$, Abaqus to $8.8$ and $9.45\times10^{-4}$, the element to $8.7\times10^{-4}$ with upwind gradients. The stress does not converge with the constant added diffusion: on $16^3$ it is about a quarter of the converged value ($1.3$ to $1.7$ against $5.9$ to $6.3\times10^{-4}$~Pa from the two second-order schemes of two codes), and on $96^3$ it is still rising. Treatments 2 to 4 give the same stress within a few per cent from $32^3$ on (Python TVD $5.90$ and $5.98\times10^{-4}$~Pa on $32^3$ and $64^3$, Abaqus SUPG $6.27$ and $6.22\times10^{-4}$~Pa on $32^3$ and $48^3$, GCI $1.2\,\%$). The element with upwind gradients gives $2.14$, $4.56$ and $5.15\times10^{-4}$~Pa on $8^3$, $16^3$ and $24^3$, extrapolated $5.8\times10^{-4}$~Pa (GCI $16.5\,\%$), within the band of the second-order schemes; its MLS operators over 30 neighbours give it a resolution of about one element, not one Gauss-point spacing (with the same added diffusion, the element on $16^3$ and Python on $32^3$ differ by $1\,\%$ in $\alpha-1$ and $5\,\%$ in the stress).

\begin{figure}[tbp]
  \centering
  \includegraphics[width=0.7\linewidth]{\figs fig_convergence.pdf}
  \caption{Separate example with antibiotic: seed means of (a)~$\alpha-1$ and (b)~the von Mises stress at $T^*=1$ against the grid, one line per scheme and code (dashed: constant added diffusion $r h/4$; solid and dotted: the second-order and the stabilised explicit forms). Grey band: the finest grids of the two second-order schemes. These runs start newly reached points from the seed composition (the alternative rule of Section~\ref{sec:newpoints}) at a field step of $0.00125$.}
  \label{fig:convergence}
\end{figure}

\subsection{Time step, regularisation and stability}
\label{app:timestep}

In the element codes the point model advances in every increment ($0.00125$ in $T^*$; Abaqus: observed order $1.1$ in time, error against the extrapolated value $0.8\,\%$ in $\alpha-1$ and $2.2\,\%$ in the stress). In the Python code the point model and $\lv$ are updated once per outer step $\Delta T^*$, inside which $\phiD$, $\alpha$ and the fields take explicit RK4 steps. The composition and $\alpha-1$ change by less than $3$ and $0.3\,\%$ between outer steps of $0.02$ and $0.0025$; the stress does not: on $64^3$ the separate example with antibiotic gives $5.68$, $5.82$, $5.94$ and $6.00\times10^{-4}$~Pa at $\Delta T^*=0.025$, $0.01$, $0.005$ and $0.0025$, the mixed one $3.23$, $3.09$, $2.90$ and $2.82\times10^{-4}$~Pa. The sensitivity comes from the front direction: after the antibiotic onset $\lv$ is small, $\nabla c$ is weak and $\nc$ responds to small changes of the consumption (a regularisation $\varepsilon^2=10^{-4}$ instead of $10^{-6}$ changes the mixed stress by $30\,\%$, $10^{-8}$ by $0.1\,\%$); the inner step, a trapezoidal update of $\lv$, consumption by newly reached points at once and the one-increment lag of $\lv$ each change the stress by less than $1\,\%$. The grid series are therefore run at $\Delta T^*=0.0025$. The explicit inner step is limited by $\beta\Delta t/h^2\le0.15$ and the Courant number of the front $Co\le0.4$; on $192^3$ both limits are close at the same time and the step diverged, which the combined condition $12D+2Co\le2.2$ removes without changing the step on $32^3$ to $128^3$.

\subsection{Grid convergence table}
\label{app:gci}

\begin{table}[htbp]
  \centering
  \caption{Grid convergence (Python, TVD, outer step $0.0025$): seed means at $T^*=1$, observed order $p$, extrapolated value and GCI \citep{Celik2008GCI}. Mixed: equal amounts unless the start fractions $(\phi_1,\phi_2)$ are given; antib.: with antibiotic; --: not yet in the asymptotic range.}
  \label{tab:gci}
  \small
  \setlength{\tabcolsep}{3.5pt}
  \begin{tabular}{llcccccc}
    \toprule
    example & quantity & $32^3$ & $64^3$ & $128^3$ & $p$ & extrapolated & GCI \\
    \midrule
    separate, antib. & $\alpha-1$ [$10^{-4}$] & $8.453$ & $8.931$ & $9.112$ & $1.4$ & $9.22$ & $1.5\,\%$ \\
    separate, antib. & von Mises [$10^{-4}$~Pa] & $5.523$ & $6.000$ & $6.284$ & $0.75$ & $6.70$ & $8.3\,\%$ \\
    separate & $\alpha-1$ [$10^{-4}$] & $8.729$ & $9.014$ & $9.110$ & $1.6$ & $9.16$ & $0.7\,\%$ \\
    separate & von Mises [$10^{-4}$~Pa] & $7.440$ & $6.730$ & $6.585$ & $2.3$ & $6.55$ & $0.7\,\%$ \\
    mixed, antib. & $\alpha-1$ [$10^{-4}$] & $7.42$ & $8.04$ & $8.32$ & $1.2$ & $8.54$ & $3.3\,\%$ \\
    mixed, antib. & von Mises [$10^{-4}$~Pa] & $2.00$ & $2.82$ & $3.30$ & $0.76$ & $3.99$ & $26\,\%$ \\
    mixed & $\alpha-1$ [$10^{-4}$] & $7.46$ & $8.08$ & $8.35$ & $1.2$ & $8.56$ & $3.1\,\%$ \\
    mixed & von Mises [$10^{-4}$~Pa] & $2.13$ & $3.08$ & $3.65$ & $0.74$ & $4.50$ & $29\,\%$ \\
    mixed $(0.25,0.30)$, antib. & $\alpha-1$ [$10^{-4}$] & $7.42$ & $8.04$ & $8.31$ & $1.2$ & $8.51$ & $3.1\,\%$ \\
    mixed $(0.25,0.30)$, antib. & von Mises [$10^{-4}$~Pa] & $2.01$ & $2.96$ & $4.01$ & -- & -- & -- \\
    \bottomrule
  \end{tabular}
\end{table}

\section*{Supplementary material}
The $L_2$ error plot of the manufactured solution, the discretisation table of the three codes, the time-step series of the neighboured element and of the Galerkin code, the stability analysis of the explicit element (Bloch analysis of the stencil, step limits $\beta\Delta t/h^2\le0.12$ to $0.16$), the sensitivities to $\phi_{\min}$, $d$, $\varepsilon$, $f$ and the start rule, the time series of the three antibiotic schedules (Section~\ref{sec:schedule}), the test of the mirrored operators of the element, and the animation of the maze will be provided as supplementary material.

{\small
\setlength{\bibsep}{2pt}
\bibliographystyle{plainnat}
\bibliography{refs}
}

\end{document}
